\section{Introduction}

Finite-context latent world models make visual control tractable by compressing
history into a short latent context.  That same compression makes memory claims
ambiguous.  A model can produce locally plausible next features while losing an
old cue that a later decision requires; a recurrent or state-space side state
can store information internally while the frozen host never exposes it; a
linear readout can recover a symbol without showing that an executed action
would change.  These are different failure modes, but they are often collapsed
into one question: does the model have memory?

We study memory as an interface property of a frozen latent host.  The host is
not retrained.  A compact causal side state observes the legal
observation/action stream, stores evidence that has left the raw context, and
writes only through the frozen host interface.  The audit is fail-closed:
stronger claims are attempted only after earlier gates pass.  The resulting
claim ladder has four rungs:
\begin{enumerate}
  \setlength\itemsep{1pt}
  \item \emph{demand}: the target cannot be solved from the legal endpoint,
        action, or state shortcuts;
  \item \emph{retention}: the side state causally carries the old evidence
        after context expiry;
  \item \emph{host exposure}: the frozen predictor exposes that evidence in
        its output rather than hiding it inside the side state; and
  \item \emph{use}: a fixed consumer can convert the exposed state into an
        executed decision.
\end{enumerate}
The paper is therefore not a claim that one carrier architecture wins.  It is a
claim that persistent memory must be audited as a target--carrier--host--
consumer chain.

\begin{figure}[!t]
  \centering
  \includegraphics[width=\linewidth]{fig_b_architecture.pdf}
  \caption{Masked-evidence memory interface.  The top strip shows native-ratio
  rendered inputs: cue-only counterfactuals share trajectory, actions, and
  endpoint, while the legal context no longer contains the old cue.  The model
  schematic follows a frozen encoder/host path with a causal evidence sidecar
  that writes through a residual host interface before exposure and use.}
  \label{fig:architecture}
\end{figure}

\paragraph{Contributions.}
This study makes four scoped contributions.  First, it turns memory
evaluation into a controlled interface audit with explicit demand, retention,
exposure, and use gates.  Second, it validates positive downstream-use cases on
PushT and PointMaze while keeping the consumer boundary explicit.  Third, it
uses Wall and OGBench to prevent overclaiming: Wall isolates a storage-versus-
exposure bottleneck, and OGBench shows that the admission protocol ports across
seven environment families while a repaired PointMaze/Cube/Puzzle/Scene
fixed-controller sweep tests whether the selected memory target changes native
environment success without claiming native OGBench planning.  Fourth, it
instantiates the audit as a
non-manual self-supervised memory objective: salient target frames and patch
sets are mined without cue labels, cue crops, or key-frame annotations, and the
slot memory predicts fixed generic color/texture latents.  This objective
passes all 30 registered OGBench PointMaze-large/Cube-single cells over ages
4/8/15 and five seeds.

\section{Audit design}

\paragraph{Frozen hosts and sidecar states.}
Each study starts from a locked finite-context latent host, either DINO-WM
\citep{r_dinowm} or LeWM \citep{r_lewm}.  The encoder and predictor remain
fixed.  A sidecar state receives the same legal stream as the host and is
evaluated under matched conditions.  The \emph{full} condition receives the old
stream.  The \emph{reset} condition reinitializes memory at the legal context
boundary.  The \emph{no-state} condition evaluates the frozen host without a
persistent side state.  These controls do not estimate an abstract scalar
``amount of memory''; they localize where the interface stops carrying old
evidence.

\paragraph{Masked evidence, not declared key frames.}
The sidecar is not given a symbolic instruction that a particular past frame
will matter.  Instead, the task constructs cue-only counterfactuals: the old
evidence changes while the later endpoint, action stream, and other legal
signals stay fixed.  The target can therefore be solved only if the side state
keeps the old evidence.  In feature-host settings the sidecar can be evaluated
both as a \emph{carrier prior} and through the \emph{host output}; separating
these two views is essential, because an internally readable carrier is not
automatically a usable world-model output.

\paragraph{Autonomous salient patch-set JEPA.}
The final OGBench method removes the remaining manual cue-feature handoff.  It
samples a causal rendered stream, mines a salient target frame and a set of
high-saliency patches without using cue labels or cue layout, masks only that
target view, and asks a compact slot state to predict stop-gradient
color/texture patch latents.  Labels are never used in the learning objective;
they enter only after training through the same full/reset/recent-only audit.
This design is intentionally not a KV cache, a recurrent hidden state, or a
manually declared key-frame mechanism.  The model must learn which evidence is
worth carrying because the generic patch-set objective repeatedly withholds
salient old observations from the legal context.

\paragraph{Admission before retention.}
Every memory claim starts with an admission deck.  Cue-time readout must verify
that the cue is visible when it appears.  Shortcut checks must show that
endpoint images, action-only traces, and state-only traces do not solve the
label.  Counterfactual masks must verify that post-cue frames do not leak the
answer.  Only after these checks pass do we run retention, exposure, or
executed-use tests.  This staging is conservative, but it is what makes a
negative or mixed result interpretable rather than merely disappointing.

\section{Environments and claim boundaries}

Figure~\ref{fig:claim-ledger} is the main evidence ledger.  The rows are not
pooled into a single score because each environment is serving a different
purpose.  PushT tests whether a visual-binding cue can influence an executed
decision on two frozen host families.  PointMaze tests whether retained goal
evidence can drive a navigation consumer.  Wall is included specifically to
expose a bottleneck where carrier storage survives but host exposure fails.
OGBench supplies breadth for the cue-admission machinery across navigation,
locomotion, manipulation, scene, and puzzle families.  Its repaired native-use
stage is reported as fixed-controller success for audited
PointMaze/Cube/Puzzle/Scene controllers, not as native world-model planning or
autonomous memory discovery.

\begin{figure}[!t]
  \centering
  \includegraphics[width=\linewidth]{fig_b_claim_ledger.pdf}
  \caption{Task-level claim ledger.  The paper avoids a pooled memory score:
  each row answers a different interface question and carries its own claim
  boundary.  A check mark means that the corresponding rung is supported by the
  current artifact; a dash means the rung is intentionally not claimed.}
  \label{fig:claim-ledger}
\end{figure}

This ledger is the central presentation choice.  A conventional results table
would invite a misleading reading in which environments are interchangeable
and carrier names can be ranked globally.  The audit instead asks where the
old evidence enters, where it disappears, and whether disappearance occurs
before or after the frozen host.

\section{Positive use cases}

The strongest evidence in this study is not a probe, but executed use.
Figure~\ref{fig:execution} compares host-readability and executed success under
fixed consumers.  The main observation is qualitative: when the host-exposed
state is reliable, a simple consumer can turn old evidence into control; when
the controls remove pre-context evidence, execution drops to the expected
uninformed regime.

\begin{figure}[!t]
  \centering
  \includegraphics[width=0.95\linewidth]{fig_b_execution.pdf}
  \caption{Readability transfers to execution when the consumer is fixed.
  PushT shows a checkpointed visual-binding use gate on DINO-WM and LeWM.
  PointMaze shows the clearest large execution separation: memory-conditioned
  goal selection succeeds where the no-state consumer behaves like an
  uninformed policy.}
  \label{fig:execution}
\end{figure}

\paragraph{PushT: use is not host-family-specific.}
The PushT binding deck asks whether a transient visual token can be retained
until a later action choice.  The result passes on both DINO-WM and LeWM frozen
hosts under the same style of full/reset/no-state comparison.  This matters
because a single-host success could be dismissed as a quirk of one feature
space.  The evidence here is narrower and cleaner: two host families can expose
the same kind of old visual binding strongly enough for a fixed consumer to
act on it.  The claim remains bounded because the banks are host-native rather
than a fully row-matched simulator comparison.

\paragraph{PointMaze: retained state can become navigation.}
PointMaze is the clearest use example because the old cue selects among
waypoint goals and the downstream controller is fixed.  The memory sidecar
does not receive a privileged goal label at decision time; it must carry the
transient mini-map evidence until the consumer reads the host output.  The
large separation between full memory and no-state execution shows that the
retained evidence is not merely decodable; it changes the chosen path under a
pinned controller.  At the same time, this is not native DINO-WM planning.  It
is evidence that a fixed external consumer can use exposed memory.

\section{Failure localization: storage is not exposure}

Wall is the most important negative result.  It was added to answer a
question that positive execution decks cannot answer: if a model fails at long
evidence age, did the side state forget, or did the frozen host stop exposing
what the side state still contains?  Figure~\ref{fig:wall} shows that these
are separable.  The carrier-prior readout remains high across ages, while the
host-output readout decays toward the shortcut regime.  The stricter
PCA-logistic confirmation also reverses an overly optimistic screening
interpretation: Wall should be written as a bottleneck diagnosis, not a
success case.

\begin{figure}[!t]
  \centering
  \includegraphics[width=0.95\linewidth]{fig_b_wall_bottleneck.pdf}
  \caption{Wall localizes the bottleneck.  Dashed curves in panel A are
  carrier-prior readouts; solid curves are host-output readouts.  The widening
  gap shows that retaining a cue internally is not the same as exposing it
  through the frozen predictor.  Panel C shows why the paper uses the stricter
  PCA-logistic confirmation instead of the earlier ridge screen.}
  \label{fig:wall}
\end{figure}

This result is useful precisely because it is not flattering.  It rules out a
simple story in which the proposed sidecar is a universally better memory
module.  Instead, the evidence points to a host-interface failure: old
information can survive in the carrier while the frozen predictor maps the
legal context to an output that no longer carries the needed variable.  A
paper that reports only carrier-prior accuracy would miss this failure.

\section{Protocol breadth: OGBench admission and native-use checks}

OGBench \citep{r_ogbench} is used as an admission bridge.  We construct
cue-only counterfactual decks for seven render-ready environments: PointMaze,
AntMaze, HumanoidMaze, Cube-single, Cube-double, Scene, and Puzzle.  The
purpose is not to claim that an OGBench world model remembers.  The purpose is
to test whether the same cue protocol can create valid old-evidence targets
without shortcut leakage across qualitatively different visual domains.  We then
run a deliberately bounded feature-host follow-up, then convert the selected
memory target into native environment success wherever a local audited
controller exists.  The repaired all-env native-use sweep scores 24/36 env-age
rows across non-teleport PointMaze, Cube-single/double/triple, Puzzle-3x3, and
Scene.  PointMaze uses a BFS waypoint controller in the native point-mass
action space.  Cube variants use the OGBench CubeMarkovOracle with live target
metadata and multi-block ordering.  Puzzle uses a lights-out solve with the
button oracle, and Scene composes the button, drawer, window, and cube Markov
oracles with dependency-aware sequencing.  PointMaze-teleport is left unscored
because its oracle-label controller gate reaches only 0.608--0.658 below the
0.850 admission threshold; AntMaze and HumanoidMaze are left unscored because
this repository has no audited low-level controller for them.  This tests
whether recovered memory changes final task success; it does not test
autonomous cue discovery or native OGBench planning.

\begin{figure}[!t]
  \centering
  \includegraphics[width=\linewidth]{fig_b_ogbench_admission.pdf}
  \caption{OGBench breadth.  Cue labels are readable at cue time, while
  endpoint, action, and state shortcuts remain at the uninformed baseline;
  counterfactual masks pass in all seven decks.  The execution panel summarizes
  the repaired fixed-controller native-use sweep: full memory converts the
  selected target into native success on the audited PointMaze/Cube/Puzzle/Scene
  rows, while recent-only and random target selection remain near chance.  This
  is fixed-controller use evidence, not native planner evidence.}
  \label{fig:ogbench}
\end{figure}

The OGBench native-use result supplies the straightforward performance number
missing from a pure readout study.  Across the 24 scored env-age rows, full
memory reaches mean native success 0.952, with worst scored success 0.828 and
mean full-minus-control lifts 0.714 over recent-only and 0.708 over random
selection.  The scientific value is not that the fixed controllers are new;
they are deliberately simple.  The value is that the same memory-selected
label that is readable at the host interface also determines whether the
environment reports success.  This fixed-controller stage still supplies the
old cue feature to the sidecar; it is therefore evidence of use, not evidence
of autonomous evidence discovery.

\paragraph{Removing the manual cue-feature handoff.}
We therefore run a final non-manual OGBench stage.  The learning rule sees only
the rendered causal stream, action tokens, and time tokens.  It mines salient
target frames and patch sets automatically, masks the selected target view,
and predicts fixed generic color/texture patch latents from a compact slot
bottleneck.  The target is not a cue crop, the model is not told which moment
is key, and the training loss never uses cue labels.  The labels are used only
after training to audit whether the learned slot state retains the old cue
under full/reset/recent-only contrasts.

\begin{figure}[!t]
  \centering
  \includegraphics[width=\linewidth]{fig_b_nonmanual_final.pdf}
  \caption{Final non-manual OGBench result.  The salient patch-set JEPA
  objective passes all 30 registered cells over two environments, three cue
  ages, and five seeds.  Full-memory readout remains high at age 15, while
  reset and recent-only controls stay near the four-way uninformed baseline.}
  \label{fig:nonmanual-final}
\end{figure}

\begin{table}[!tbp]
\centering
\footnotesize
\setlength{\tabcolsep}{4.5pt}
\renewcommand{\arraystretch}{1.08}
\caption{Final non-manual OGBench summary.  Gate: full balanced accuracy
at least 0.750 and reset/recent-only controls at most 0.350 for every seed.}
\label{tab:nonmanual-final}
\begin{tabular}{lcccccc}
\toprule
Deck & Age & Seeds & Full & Reset & Recent-only & Status\\
\midrule
Cube-single & 4  & 5/5 & 1.000 & 0.240 & 0.257 & \Pass\\
Cube-single & 8  & 5/5 & 0.990 & 0.223 & 0.254 & \Pass\\
Cube-single & 15 & 5/5 & 0.921 & 0.247 & 0.246 & \Pass\\
PointMaze-large & 4  & 5/5 & 1.000 & 0.251 & 0.255 & \Pass\\
PointMaze-large & 8  & 5/5 & 0.994 & 0.266 & 0.250 & \Pass\\
PointMaze-large & 15 & 5/5 & 0.983 & 0.251 & 0.252 & \Pass\\
\bottomrule
\end{tabular}
\end{table}

\paragraph{Breadth stress tests.}
The two-deck result above is the current clean paper-level non-manual claim.
The broader PointMaze/Cube/Scene/Puzzle/locomotion sweep is treated as a
stress test rather than an automatic headline.  Its generated status is
recorded in Appendix~\ref{app:nonmanual-breadth-auto}; rows that require
large-bank confirmation are reported as such instead of being folded into a
single pooled score.

\section{Cross-study interpretation}

\paragraph{The interface chain, not the carrier name, is the unit of evidence.}
The same sidecar can look strong under a carrier-prior readout and weak under
host-output exposure.  Conversely, a modest readout can still matter if a
consumer is sensitive to the recovered variable.  This is why the paper avoids
phrases such as ``fixed-trust solves memory'' or ``state-space is the memory
winner.''  The defensible statement is conditional: for a given target, host,
sidecar, and consumer, the old evidence either survives the chain or it does
not.

\paragraph{Negative results sharpen the method.}
The Wall result is not dead weight; it is the diagnostic that makes the
positive PushT and PointMaze results credible.  Without Wall, the paper could
be read as another architecture-comparison story.  With Wall, the paper shows
why architecture comparisons are brittle: a carrier can preserve the symbol
while the host interface erases it.  That is the phenomenon reviewers should
remember.

\paragraph{Execution should be interpreted as consumer-specific use.}
The fixed consumers are deliberately simple and shared within each deck.  They
prove that exposed memory can change control, not that the native planner has
learned to query memory.  This distinction is not a weakness to hide.  It is
the correct boundary between an interface audit and an end-to-end planning
paper.

\paragraph{Breadth should not outrun evidence.}
The OGBench decks show that the admission stage generalizes beyond the original
hosts and tasks, and the expanded fixed-controller use stage shows that
recovered memory can change native environment success across audited
PointMaze/Cube/Puzzle/Scene rows.  Treating this as native planner memory
would inflate the paper but weaken its rigor.  The
non-manual patch-set JEPA stage takes the next step by removing the explicit
cue-feature handoff on two core OGBench environments.  The right boundary is
therefore: autonomous old-evidence retention is supported for
PointMaze-large/Cube-single render-stage memory, while native OGBench
world-model planning and broad OGBench-wide autonomous generalization remain
unclaimed.

\section{Diagnostic consequences}

The audit is useful only if its outcomes imply different next actions.
Table~\ref{tab:diagnosis} summarizes the signatures exposed by the current
suite.  The same endpoint score can have different causes depending on the
carrier prior, reset, no-state, host output, and execution views.

\begin{table}[!tbp]
\centering
\footnotesize
\setlength{\tabcolsep}{3.8pt}
\renewcommand{\arraystretch}{1.12}
\caption{How to read audit signatures.  The entries are written as diagnoses,
not as leaderboard ranks.}
\label{tab:diagnosis}
\begin{tabularx}{\linewidth}{lY Y}
\toprule
Observed signature & Interpretation & Correct intervention target\\
\midrule
Prior weak, host output weak, controls at chance
  & The side state did not preserve the old evidence in a useful form.
  & Memory objective, write dynamics, capacity, or training horizon.\\
Prior strong, host output weak
  & The carrier stored evidence, but the frozen fusion/predictor path did not
    expose it.
  & Exposure adapter, fusion point, residual coordinates, or host-aware
    auxiliary objective.\\
Host output readable, execution unchanged
  & The representation is accessible to a probe but not used by the declared
    consumer.
  & Consumer interface, action grounding, or planner integration.\\
Full beats no state but not reset
  & The sidecar helps the endpoint, but dependence on pre-context state is not
    isolated under the reset intervention.
  & Report as sensitivity, not persistent retention; redesign reset or task.\\
Admission passes but no host stage exists
  & The task is clean, not yet a memory experiment.
  & Add a frozen host, carrier grid, reset/no-state controls, and use deck.\\
\bottomrule
\end{tabularx}
\end{table}

This decision table clarifies why Wall is valuable.  A positive-only paper
would treat Wall as a failed environment and move on.  A memory audit treats it
as evidence about the interface.  The prior is strong enough to rule out simple
forgetting inside the fixed-trust carrier; the host output is weak enough to
rule out successful exposure.  That points to a concrete repair hypothesis: a
self-supervised memory method should not only store masked evidence, but also
train the host-facing projection to keep that evidence in coordinates that
survive frozen predictor mixing.

\section{Environment roles and why more rows are not automatically better}

The current environments are not interchangeable.  PushT, PointMaze, Wall, and
OGBench each stress a different part of the ladder.  Table~\ref{tab:envroles}
states the role of each environment in the argument and keeps the paper from
reading like an experiment log.

\begin{table}[!tbp]
\centering
\footnotesize
\setlength{\tabcolsep}{3.5pt}
\renewcommand{\arraystretch}{1.10}
\caption{Environment roles in the paper narrative.}
\label{tab:envroles}
\begin{tabularx}{\linewidth}{lY Y Y}
\toprule
Environment & What it contributes & What it does not prove & Why it belongs in main text\\
\midrule
PushT & Visual binding under two frozen host families with an executed-use gate.
  & A row-matched causal comparison between DINO-WM and LeWM.
  & Shows the use gate is not host-family specific.\\
PointMaze & Goal memory that survives host readout and changes navigation under
  a fixed controller.
  & Native model planning with memory.
  & Provides the strongest positive behavioral link.\\
Wall & A controlled negative case where storage and exposure separate.
  & Broad success of the proposed carrier.
  & Makes the diagnostic thesis falsifiable and prevents overclaiming.\\
OGBench & Seven clean cue decks, plus repaired PointMaze/Cube/Puzzle/Scene
  native success under fixed controllers after memory-selected target recovery,
  plus a non-manual salient patch-set JEPA stage on PointMaze-large/Cube-single.
  & Native OGBench world-model planning or broad OGBench-wide autonomous
    generalization.
  & Gives a reviewer-facing final-performance check and a clean autonomous
    memory-learning result beyond linear readout.\\
\bottomrule
\end{tabularx}
\end{table}

This structure is stronger than adding many partially completed environments.
The paper now spans positive use, exposure failure, protocol breadth, and
fixed-controller native success.  The missing piece is not another admission
deck; it is an autonomous evidence-discovery stage and audited locomotion
controllers.  AntMaze remains the most useful next execution target for
locomotion generalization, but it requires a real ant controller or policy
before it can be promoted beyond feature-host capacity.

\section{Implications for memory-module design}

The results suggest a design principle for future modules: self-supervised
memory should optimize an evidence interface, not merely a hidden state.  A
module that reconstructs or predicts masked old evidence internally can still
fail if its host-facing projection is not trained to expose that evidence.  A
module that exposes evidence can still fail if the consumer reads the wrong
coordinates.  Therefore the objective should couple three pressures: retain
masked evidence without an explicit key-frame label, expose it through the
frozen host interface, and preserve downstream separability under reset and
no-state controls.

This principle is compatible with JEPA-style training.  The sidecar need not
decode pixels or be told which observation is important for the final decision.
Instead, the final OGBench stage mines salient target frames and patch sets
without labels, predicts fixed generic color/texture latents, and audits the
result only after training.  The important lesson from Wall is that this target
should remain host-facing: optimizing only the carrier prior risks building a
memory that exists in the side state but is invisible after frozen predictor
mixing.  The important lesson from PointMaze and Cube is that a simple
self-supervised patch-set target can be sufficient when the retained
coordinates preserve the old evidence under reset and recent-only controls.

For paper framing, this means the study should not sell fixed-trust as the
new winning architecture.  The stronger claim is that the audit identifies the
right unit of design---the evidence interface---and that salient patch-set
JEPA is one concrete non-manual instantiation that passes a registered
PointMaze/Cube grid.  The remaining claim boundary is breadth and native
planning: the method has not yet been certified across every admitted OGBench
deck or inside an end-to-end world-model planner.

\section{Related work and positioning}

Recent work adds memory to world models through recurrent or state-space
agents, transformer memory, explicit retrieval, spatial/geometric memory, and
world-action memory systems
\citep{r_r2i,r_onmemory,r_s4wm,r_memworld,r_memorywam,r_dimwam,r_persist,
r_statecore}.  These papers primarily ask how to build stronger memory-enabled
agents or simulators, often evaluating through return, rollout consistency,
video reconstruction, or task success.  This paper asks a different question:
under a frozen finite-context latent host, which interface rung actually
carries the old evidence?

The closest host background is DINO-WM \citep{r_dinowm} and LeWorldModel
\citep{r_lewm}.  We freeze these hosts and treat memory as a side-interface
audit rather than a new host training recipe.  This positioning is also
consistent with broader cautions about probing and planning: decodability does
not by itself establish task use \citep{r_belinkov,r_ravich}, and predictive
feature quality does not guarantee downstream planning utility \citep{r_rcaux}.

\section{Limitations}

The main limitations are claim-boundary limitations.  First, the paper is not
a universal architecture win.  Wall shows that an internally strong carrier can
fail at host exposure, and the claim ledger should keep this visible.  Second,
the strongest execution results use fixed external consumers and pinned
controllers, not native planners; PointMaze-teleport fails the controller gate
and Ant/Humanoid locomotion remains unscored without a local low-level
controller.  Third, the non-manual patch-set JEPA result is currently certified
on PointMaze-large and Cube-single, not across every admitted OGBench deck.
Fourth, OGBench still does not include a native world-model planner that
autonomously queries the learned memory.  Fifth, the PushT cross-host result establishes that the use gate can pass on
two host families, but it is not a row-matched causal comparison between those
families.  These limitations are not cosmetic; they define what the audit can
and cannot certify.

\section{Conclusion}

The evidence supports a concise ICLR story: persistent memory in
finite-context latent world models is an interface property.  The side state
may retain old evidence, the frozen host may expose or suppress it, and a
consumer can use it only when the exposed representation remains reliable.
PushT and PointMaze provide positive executed-use cases; Wall identifies the
storage-versus-exposure bottleneck; OGBench shows that valid old-evidence
admission decks can be carried through to repaired 24/36 env-age
fixed-controller native-use success and that a non-manual salient patch-set
JEPA objective can retain old evidence on PointMaze-large and Cube-single
without cue labels, cue crops, or key-frame annotations.  The
resulting claim is deliberately scoped: the audit can certify where evidence
survives the interface chain, but it does not turn any single carrier name into
a universal memory solution.
